% =========================================================
% Paper revisions after Phase A (rolling-origin) and Phase D
% (paired statistical testing). Drop-in replacement blocks for
% report/paper.tex. The original paper.tex is left untouched.
%
% Section labels match paper.tex so cross-references continue to
% resolve.
% =========================================================

% ---------------------------------------------------------------------
% 1. ABSTRACT (replace existing abstract environment)
% ---------------------------------------------------------------------
\begin{abstract}
Bitcoin-OTC was a peer-rated trust marketplace whose users rated each
other on a $[-10, +10]$ scale, producing a signed, weighted, directed,
timestamped social network. We study early warning of \emph{new
negative} trust edges: at the end of month $t$, can we rank ordered
pairs of users by their risk of receiving a new negative directed
rating in month $t{+}1$, using only information available at time $t$?
We build an interpretable, leakage-controlled pipeline with cumulative
monthly snapshots and a three-way feature ablation
(structural; structural + temporal; structural + temporal + community-aware)
under Logistic Regression and Random Forest, with Louvain communities
extracted per snapshot. We evaluate under a rolling-origin protocol
over 25 disjoint single-month test folds covering 425 total positive
events. Adding temporal features yields a statistically significant
ROC-AUC improvement under both models (paired Wilcoxon
$p \leq 1.4\!\times\!10^{-4}$ after Holm-Bonferroni correction on
four pre-declared primary tests; matched-pairs rank-biserial
$r_{\mathrm{rb}} \approx +0.8$), and a significant Random Forest PR-AUC
improvement ($p = 6.3\!\times\!10^{-4}$, $r_{\mathrm{rb}} = +0.74$).
Adding community-aware features on top of temporal features yields a
small, directionally consistent Random Forest PR-AUC gain (19 of 25
folds positive; sign test $p = 0.015$; Wilcoxon $p = 0.048$, marginal
under the pre-declared threshold). Community features do not
significantly improve ROC-AUC over the temporal baseline ($p = 0.096$,
10/25 folds positive), and yield no detectable gain under Logistic
Regression on any metric. Precision@$k$ proves unstable under the
rolling protocol -- per-fold median P@50 is zero -- and is reported
as a supplementary metric. We frame temporal lift as the robust
headline finding and community lift as a modest, PR-oriented
refinement under Random Forest.
\end{abstract}

% ---------------------------------------------------------------------
% 2. INTRODUCTION -- contributions list (replace the existing itemize)
% ---------------------------------------------------------------------
Our contributions are:
\begin{itemize}
\item An interpretable feature-based pipeline -- classical
link-prediction, signed structural, trust-summary, temporal, and
community-aware features -- with explicit leakage controls and a
time-based candidate generation procedure.
\item A three-way feature ablation (structural; structural + temporal;
structural + temporal + community) under Logistic Regression and
Random Forest.
\item A \emph{rolling-origin evaluation} over 25 disjoint single-month
test folds (425 total positive events), with paired per-fold metric
deltas across feature sets.
\item Formal paired statistical testing on per-fold deltas: two-sided
Wilcoxon signed-rank and exact binomial sign tests, with
Holm-Bonferroni correction on a pre-declared primary family of four
tests and Benjamini-Hochberg FDR control on twenty secondary tests.
\item A descriptive analysis of community structure and weak-tie
signals, separating global feature-level evidence from a small local
top-of-ranked-list diagnostic.
\end{itemize}

The pipeline is intentionally lighter than the deep signed-graph
methods surveyed by Zhang \emph{et al.}~\cite{zhang2024}. Our goal is
interpretable analysis on a single dataset; we do not claim
state-of-the-art absolute performance.

% ---------------------------------------------------------------------
% 3. SECTION 5 -- Experimental Setup (replace whole section)
% ---------------------------------------------------------------------
\section{Experimental Setup}\label{sec:setup}

\subsection{Eligible snapshot pool}
A snapshot is \emph{eligible} for evaluation if (i) its next-month
window contains at least \code{MIN\_POS\_FOR\_SPLIT = 10} raw new
negative events, and (ii) it sits in the longest contiguous block of
eligible snapshots (this excludes a sporadic late-dataset island at
$t = 53$). The resulting eligible pool is $t \in \{13, \dots, 49\}$
(37 snapshots).

\subsection{Rolling-origin evaluation protocol (primary)}
\label{sec:setup:rolling}

Our primary evaluation is an expanding-window rolling-origin protocol
over the eligible pool. For each fold index
$k \in \{25, 26, \dots, 49\}$ (25 disjoint single-month test folds):
\begin{itemize}
\item TRAIN  = candidate pairs from snapshots $\{13, \dots, k{-}2\}$,
\item VAL    = candidate pairs from snapshot $\{k{-}1\}$ (held out; not used for tuning),
\item TEST   = candidate pairs from snapshot $\{k\}$.
\end{itemize}
We use an expanding-window TRAIN rather than a sliding window because
positive events are rare and discarding early ones would weaken
training. We use single-month TEST windows rather than two-month
windows because they maximize the number of disjoint folds (25 vs
$\sim$12), match an operational ``next-month watchlist'' deployment
cadence, and produce fold-level metric vectors directly usable for
paired non-parametric tests. The total number of positive test events
across all folds is $\sum_{k} n_{\mathrm{pos}}(k) = 425$, an
approximately $12\!\times$ increase over the single locked fold
described in Section~\ref{sec:setup:locked}.

For each fold and each (model, feature set) combination, we fit a
fresh model on the fold's TRAIN with the same $100{:}1$ training
negative subsampling as in Section~\ref{sec:method}, score the fold's
TEST, and record ROC-AUC, PR-AUC, Precision@50, and Precision@100.
All seeds are fixed; within a fold the training subsample is
identical across (model, feature set), so per-fold deltas across
feature sets are properly paired.

\subsection{Paired statistical testing on per-fold deltas}
\label{sec:setup:stats}

On the $n = 25$ per-fold deltas we compute, per (model, comparison,
metric):
\begin{itemize}
\item a two-sided paired Wilcoxon signed-rank test
(\code{zero\_method='wilcox'});
\item a two-sided exact binomial sign test on the direction counts;
\item the matched-pairs rank-biserial correlation
$r_{\mathrm{rb}} = (W_+ - W_-)/(W_+ + W_-)$ as a Wilcoxon effect
size.
\end{itemize}

We declare a small \emph{primary family} of four tests in advance --
the ones we consider central to the paper's question: RF (S+T) - S on
ROC-AUC, RF (S+T) - S on PR-AUC, RF (S+T+C) - (S+T) on PR-AUC, and
LR (S+T) - S on ROC-AUC. We control family-wise error at
$\alpha = 0.05$ on this primary set with Holm-Bonferroni. All other
(model, comparison, metric) combinations form a secondary family, on
which we control the false-discovery rate at $q = 0.05$ with
Benjamini-Hochberg. We report raw p-values, the corrected reject
decision, and the rank-biserial effect size for every test. When the
per-fold delta is exactly zero on many folds (particularly for
Precision@$k$, where the discrete probability grid produces many
ties), the Wilcoxon test becomes unreliable and we rely on the sign
test only.

\subsection{Held-out single-fold analysis (secondary)}
\label{sec:setup:locked}

For continuity with prior pipeline stages and to support the bootstrap
analysis in Section~\ref{sec:results:bootstrap}, we also retain the
last two eligible snapshots ($t \in \{48, 49\}$) as a single locked
held-out fold, with the immediately preceding two snapshots
($t \in \{46, 47\}$) as a tuning-free validation set
(Table~\ref{tab:split}). All numbers from this single fold are
reported in Section~\ref{sec:results:locked}; they are kept as a
secondary analysis. The headline numerical claims in this paper come
from the rolling-origin protocol.

\begin{table}[t]
\caption{Locked single held-out fold. Used in
Section~\ref{sec:results:locked} for the bootstrap analysis and curve
overlays; the headline numerical claims come from the rolling-origin
protocol of Section~\ref{sec:setup:rolling}.}
\label{tab:split}
\centering
\begin{tabular}{lccc}
\toprule
Split & Snapshots & Pairs & Positives \\
\midrule
TRAIN & $t = 13 \dots 45$ & 5{,}013{,}346 & 545 \\
VAL   & $t = 46, 47$      &   184{,}648   &  16 \\
TEST  & $t = 48, 49$      &   148{,}639   &  34 \\
\bottomrule
\end{tabular}
\end{table}

\subsection{Bootstrap protocol (for the locked single fold)}
\label{sec:setup:bootstrap}
For the locked single-fold analysis we attach uncertainty estimates
with a paired stratified bootstrap. For $B = 1000$ iterations, we
draw 34 positive indices with replacement and 148{,}605 negative
indices with replacement (preserving the original prevalence and
guaranteeing both labels in every resample). The same resample index
vector is reused across all three feature sets within one model, so
deltas such as
$(\mathrm{S}{+}\mathrm{T}{+}\mathrm{C}) - (\mathrm{S}{+}\mathrm{T})$
are properly paired. We add a $\mathcal{U}[0, 10^{-12}]$ jitter to RF
probabilities before top-$k$ ranking to randomize ties on the
discrete $1/200$ leaf-vote grid without perturbing non-tied rankings.
We report 95\% percentile confidence intervals and the
bootstrap-positive share $\mathrm{frac\_boot}(>0)$; this share is
\emph{not} a frequentist p-value, and we interpret a CI that overlaps
zero as ``not distinguishable from sampling noise on this test set.''
This analysis quantifies uncertainty conditional on the locked fold;
the rolling-origin protocol of Section~\ref{sec:setup:rolling}
provides the cross-fold significance picture.

\subsection{Reproducibility and artifacts}
All seeds (model fits, training subsample, Louvain partition,
bootstrap) are fixed at $42$. Intermediate parquet files and the
locked-fold test-probability matrix
(\code{results/test\_predictions.parquet}) are persisted so that
downstream analysis re-runs without re-training. Per-fold rolling-origin probabilities for the Random Forest S+T+C model are stored
under \code{results/rolling\_predictions/}. Code is organized in
\code{src/} and \code{scripts/}; the headline numbers in
Section~\ref{sec:results} come from
\code{results/rolling\_eval\_per\_fold.csv} (rolling-origin),
\code{results/rolling\_stats\_tests.csv} (paired tests), and
\code{results/test\_metrics\_from\_predictions.csv} (locked fold).

% ---------------------------------------------------------------------
% 4. SECTION 6 -- Results (replace whole section through end of
% bootstrap subsection; subsubsection notes mark figures/tables to
% keep from the previous draft verbatim)
% ---------------------------------------------------------------------
\section{Results}\label{sec:results}

We label feature sets S (structural + trust summary, 16 columns),
S$+$T (+temporal, 28 columns), and S$+$T$+$C (+community, 36 columns).

\subsection{Rolling-origin metrics}\label{sec:results:rolling}

Table~\ref{tab:rolling_summary} reports across-fold mean and standard
deviation per (model, feature set) over the $n = 25$ rolling-origin
folds (total 425 positive test events). Adding temporal features
clearly improves both ROC-AUC and PR-AUC under both models; adding
community-aware features on top of temporal moves PR-AUC slightly
upward under Random Forest while leaving Logistic Regression
essentially unchanged.

\begin{table}[t]
\caption{Rolling-origin metrics across 25 single-month test folds
(total 425 positives). Values are across-fold mean $\pm$ std.}
\label{tab:rolling_summary}
\centering
\small
\begin{tabular}{llcccc}
\toprule
F.\,set & Model & ROC-AUC & PR-AUC & P@50 & P@100 \\
\midrule
S       & lr & $0.889 \pm 0.077$ & $0.006 \pm 0.006$ & $0.012$ & $0.008$ \\
S$+$T   & lr & $0.930 \pm 0.046$ & $0.007 \pm 0.008$ & $0.014$ & $0.010$ \\
S$+$T$+$C & lr & $0.928 \pm 0.048$ & $0.007 \pm 0.009$ & $0.014$ & $0.010$ \\
\addlinespace
S       & rf & $0.881 \pm 0.073$ & $0.006 \pm 0.011$ & $0.009$ & $0.008$ \\
S$+$T   & rf & $0.924 \pm 0.054$ & $0.014 \pm 0.023$ & $0.022$ & $0.016$ \\
S$+$T$+$C & rf & $0.920 \pm 0.058$ & $0.017 \pm 0.029$ & $0.022$ & $0.016$ \\
\bottomrule
\end{tabular}
\end{table}

\subsection{Paired statistical tests}\label{sec:results:tests}

Table~\ref{tab:primary_tests} reports the four pre-declared primary
tests on the per-fold deltas, with Holm-Bonferroni decisions at
$\alpha = 0.05$. All four nulls are rejected after correction, but
the RF community PR-AUC test is marginal -- its Wilcoxon $p = 0.048$
sits at the loosest Holm threshold. The sign-test p-value
($p = 0.015$) on the same comparison is more robust to ties.
Fig.~\ref{fig:paired_deltas} visualizes the per-fold deltas for
these four comparisons.

\begin{table}[t]
\caption{Primary paired statistical tests on rolling-origin per-fold
deltas ($n = 25$; two-sided paired Wilcoxon signed-rank;
Holm-Bonferroni at $\alpha = 0.05$). $r_{\mathrm{rb}}$ is the
matched-pairs rank-biserial effect size.}
\label{tab:primary_tests}
\centering
\small
\begin{tabular}{llccc}
\toprule
Comparison & Metric & Wilcoxon $p$ & sign $p$ & $r_{\mathrm{rb}}$ \\
\midrule
RF $(\mathrm{S}{+}\mathrm{T}) - \mathrm{S}$       & ROC-AUC & $4.5{\times}10^{-5}^{\dagger}$ & $1.6{\times}10^{-4}$ & $+0.85$ \\
LR $(\mathrm{S}{+}\mathrm{T}) - \mathrm{S}$       & ROC-AUC & $1.4{\times}10^{-4}^{\dagger}$ & $9.1{\times}10^{-4}$ & $+0.81$ \\
RF $(\mathrm{S}{+}\mathrm{T}) - \mathrm{S}$       & PR-AUC  & $6.3{\times}10^{-4}^{\dagger}$ & $9.1{\times}10^{-4}$ & $+0.74$ \\
RF $(\mathrm{S}{+}\mathrm{T}{+}\mathrm{C}) - (\mathrm{S}{+}\mathrm{T})$ & PR-AUC  & $4.8{\times}10^{-2}^{\dagger}$ & $1.5{\times}10^{-2}$ & $+0.45$ \\
\bottomrule
\end{tabular}\\[2pt]
\raggedright\footnotesize${}^{\dagger}$ rejects $H_0$ after Holm-Bonferroni at $\alpha = 0.05$.
\end{table}

\begin{figure}[t]
\centering
\includegraphics[width=\columnwidth]{results/figures/rolling_paired_deltas.png}
\caption{Per-fold paired deltas for the four primary comparisons
(blue: positive fold; red: negative fold).}
\label{fig:paired_deltas}
\end{figure}

Among secondary tests (Benjamini-Hochberg FDR $q = 0.05$): the full
$\mathrm{S} \to \mathrm{S}{+}\mathrm{T}{+}\mathrm{C}$ delta is
significant on both metrics and both models;
$\mathrm{LR}\,(\mathrm{S}{+}\mathrm{T}) - \mathrm{S}$ on PR-AUC is
significant despite a small absolute magnitude;
$\mathrm{RF}\,(\mathrm{S}{+}\mathrm{T}) - \mathrm{S}$ on
Precision@100 is significant. No community-vs-temporal ROC-AUC test
reaches significance: RF (S+T+C) - (S+T) ROC-AUC has $p = 0.096$ with
only 10/25 folds positive, consistent with saturation; LR (S+T+C) -
(S+T) ROC-AUC has $p = 0.47$. LR community-vs-temporal PR-AUC is
also null ($p = 0.35$). Most Precision@$k$ tests either fail to
reach significance after correction or are flagged unreliable
because of tied zeros.

\subsection{Temporal ablation}\label{sec:results:temporal}

The temporal-over-structural ROC-AUC delta is positive in 22 of 25
RF folds (median $+0.031$, mean $+0.044$) and in 21 of 25 LR folds
(median $+0.036$, mean $+0.041$). Both pass Holm correction with
very small p-values and large effect sizes
($r_{\mathrm{rb}} \approx +0.85$ and $+0.81$).

The Random Forest temporal PR-AUC delta is also significant after
Holm correction (Wilcoxon $p = 6.3\!\times\!10^{-4}$, sign-test
$p = 9.1\!\times\!10^{-4}$, $r_{\mathrm{rb}} = +0.74$), with 21 of
25 folds positive. The Logistic Regression temporal PR-AUC delta is
significant under the secondary BH correction
($p = 3.4\!\times\!10^{-3}$), but its mean magnitude ($+0.002$) is
small in absolute terms.

The temporal feature family is internally redundant by design:
recent-window counts and exponentially decayed aggregates carry
overlapping information, with within-pair correlations $\geq 0.9$
for some pairs. We retain both for transparency.

\subsection{Community-aware ablation}\label{sec:results:community}

\textbf{Community PR-AUC under Random Forest is small but
directionally consistent.} The (S+T+C) - (S+T) PR-AUC delta is
positive in 19 of 25 folds, with median $+0.001$ and mean $+0.003$.
The Wilcoxon test sits right at the Holm threshold
($p = 0.048$); the sign test is more robust ($p = 0.015$). The
rank-biserial effect size $r_{\mathrm{rb}} = +0.45$ is in the medium
range. We describe this gain as small but directionally consistent
across folds, with marginal statistical support.

\textbf{Community does not improve ROC-AUC over temporal.} The
(S+T+C) - (S+T) ROC-AUC delta is positive in only 10 of 25 RF folds
($p = 0.096$, $r_{\mathrm{rb}} = -0.39$) and 12 of 25 LR folds
($p = 0.47$). The Random Forest direction is mildly negative on
average, consistent with ROC-AUC saturating at the $+$T level rather
than with community features being harmful.

\textbf{Logistic Regression shows no community gain on any metric.}
This is consistent with the community signal being non-linear and
inaccessible to a linear model on these features.

At the descriptive level, the data are not consistent with the
naive ``distrust at community boundaries'' reading. Aggregating
across all rolling-fold training blocks, same-community pairs have
an empirical positive rate of $0.0140\%$ ($215 / 1{,}540{,}969$)
while cross-community pairs are at $0.0095\%$ ($330 / 3{,}472{,}377$);
same-community pairs are slightly \emph{more} likely to receive a
new negative rating in the next window (Fig.~\ref{fig:samecross}).

\begin{figure}[t]
\centering
\includegraphics[width=0.85\columnwidth]{results/figures/same_vs_cross_positive_rate.png}
\caption{Aggregate training positive rate by community membership.
Same-community pairs are $\sim 47\%$ more likely to receive a new
negative rating than cross-community pairs. The naive
``cross-community distrust'' hypothesis is not supported.}
\label{fig:samecross}
\end{figure}

In the fitted Logistic Regression on S$+$T$+$C,
\code{neighborhood\_overlap} is the strongest single community
feature (standardized coefficient $-1.36$, in the direction that
lower shared-neighborhood predicts higher new-negative probability);
it is also among the top-ranked community features in Random Forest
(Tables~\ref{tab:imp_lr_topfam} and \ref{tab:imp_rf_topfam} in
Section~\ref{sec:results:locked}). The coarse same/cross indicator
carries little weight once \code{neighborhood\_overlap}, boundary
fractions, and community sizes are present.

\textbf{Precision@$k$ is unstable under rolling.} Per-fold median
P@50 is zero for every (model, feature set), and median P@100 is at
most $0.01$ for the best (model, feature set) combination. Most
folds return no true positives in the top 50. We therefore demote
Precision@$k$ to a supplementary metric.

\subsection{Held-out single-fold analysis (secondary)}
\label{sec:results:locked}

For continuity with prior pipeline stages,
Table~\ref{tab:test_metrics} reports the locked single held-out fold
metrics. The locked fold's RF (S+T+C) PR-AUC ($0.0703$) is well
above the rolling-fold mean ($0.017$); the locked fold thus
over-represents the magnitude of the community PR-AUC effect that is
visible across all 25 folds. We retain the table, the curve
overlays, the importance rankings, and the bootstrap analysis here
because they remain informative about model behavior on a single
representative fold, but the headline numerical claims in this paper
come from the rolling-origin protocol.

\begin{table}[t]
\caption{Held-out single-fold metrics ($t \in \{48, 49\}$;
34 positives, 148{,}639 candidate pairs). Secondary to
Table~\ref{tab:rolling_summary}.}
\label{tab:test_metrics}
\centering
\begin{tabular}{llcccc}
\toprule
F.\,set & Model & ROC-AUC & PR-AUC & P@50 & P@100 \\
\midrule
S      & logreg & 0.909 & 0.0034 & 0.00 & 0.01 \\
S$+$T   & logreg & 0.949 & 0.0059 & 0.00 & 0.01 \\
S$+$T$+$C & logreg & 0.945 & 0.0058 & 0.00 & 0.01 \\
\addlinespace
S      & rf     & 0.880 & 0.0223 & 0.04 & 0.03 \\
S$+$T   & rf     & 0.950 & 0.0400 & 0.08 & 0.07 \\
S$+$T$+$C & rf     & 0.940 & 0.0703 & 0.10 & 0.06 \\
\bottomrule
\end{tabular}
\end{table}

\subsubsection{ROC and PR curves on the locked fold}
% Keep fig:rf_roc and fig:rf_pr verbatim from the previous draft.
Fig.~\ref{fig:rf_roc} and Fig.~\ref{fig:rf_pr} overlay Random Forest
ROC and PR curves for the three feature sets on the locked fold. The
S$+$T and S$+$T$+$C ROC curves nearly coincide in the operationally
relevant low-FPR region; the S$+$T$+$C PR curve sits visibly above
S$+$T between recall $\sim 0.05$ and $\sim 0.30$. Under the rolling
protocol the average PR separation is much smaller
(Table~\ref{tab:rolling_summary}); the locked fold is on the
favorable end of the per-fold distribution.

\subsubsection{Feature importance (qualitative)}
% Keep tab:imp_lr_topfam and tab:imp_rf_topfam verbatim from
% the previous draft. Their narrative role is now qualitative
% support for the rolling-origin claims.

\subsubsection{Bootstrap uncertainty on the locked fold}
\label{sec:results:bootstrap}
% Keep tab:boot_abs and tab:boot_delta verbatim from the previous
% draft. Replace the closing bullet list with this short note.
These bootstrap CIs quantify sampling uncertainty conditional on the
locked fold. They were informative when the locked fold was the only
held-out evaluation; with the rolling-origin protocol in place they
are complementary rather than primary. Note that the locked-fold
community PR-AUC effect ($+0.026$, 94\% of bootstrap resamples
positive, CI $[-0.003, +0.068]$) is consistent in direction with the
rolling result but is several times larger than the across-fold mean
($+0.003$).

% ---------------------------------------------------------------------
% 5. SECTION 7 -- Error Analysis: ONLY the opener of subsection
% "Global, feature-level evidence" changes; the two enumerated facts
% and the rest of section 7 stay verbatim.
% ---------------------------------------------------------------------
\subsection{Global, feature-level evidence}

Two descriptive facts at the population level are worth stating
separately. They are aggregated across all 25 rolling-fold training
blocks; the same direction holds for any individual fold's training
block.
% (the rest of the subsection is unchanged: the two enumerated facts
% and the intra-community-weak-tie reading paragraph stay verbatim.)

% ---------------------------------------------------------------------
% 6. SECTION 8 -- Limitations (replace whole section)
% ---------------------------------------------------------------------
\section{Limitations}\label{sec:lim}

We list the limitations explicitly so that follow-up work can
prioritize.

\begin{itemize}
\item \textbf{Per-fold positive counts are still small.} Although
the rolling-origin protocol accumulates 425 positive test events
across 25 folds, individual folds have only $\sim 17$ positives on
average, and per-fold PR-AUC has standard deviation roughly twice
its mean for Random Forest. This is why the community PR-AUC effect
sits at the margin of statistical significance even with $n = 25$.
\item \textbf{Single network, no external replication.} We do not
report results on Bitcoin-Alpha or any other signed network.
Bitcoin-Alpha is a natural transfer target and the most important
next step.
\item \textbf{Stronger baselines are missing.} Our comparisons are
between feature groups under two interpretable models. We do not
benchmark gradient-boosted trees, simple time-aware heuristics
(recency-only, victim-risk-only), or shallow graph embeddings; any
of these could absorb part of what we report as ``temporal lift''
or ``community lift.''
\item \textbf{Precision@$k$ is unstable under the rolling
protocol.} Per-fold median P@50 is zero for every (model, feature
set). We report P@$k$ only as supplementary evidence and do not use
it to support community claims.
\item \textbf{Last-rating-wins snapshot collapse.} The snapshot
graph uses the latest rating per pair; this represents current
trust judgment but discards multi-event nuance. We restore the full
event history only for temporal and trust-summary features.
\item \textbf{Bootstrap CIs apply to the locked single fold.}
Section~\ref{sec:setup:bootstrap}'s percentile CIs quantify sampling
uncertainty conditional on the locked fold and can under-cover when
the bootstrap distribution is skewed (e.g.\ for Precision@$k$ near
its hard ceiling). The rolling-origin Wilcoxon results are our
primary uncertainty story.
\item \textbf{Pre-declared primary family and multiple-testing
discipline.} The primary family of four tests and the Holm
correction were fixed before inspecting the rolling-origin results,
but the community PR-AUC primary test passes only marginally
($p = 0.048$). Re-running with a different primary family, a
different correction (e.g.\ Bonferroni on all 24 tests at
$\alpha / 24 = 0.002$), or a different metric set could flip this
single decision. We report this honestly rather than
soft-pedalling it.
\item \textbf{\code{prior\_interaction} is partly degenerate.} By
construction of the candidate set, this flag is $1$ only when $u$
previously rated $v$ positively. We retain it for transparency but
note that its fitted coefficient should be read with care.
\item \textbf{\emph{frac\_boot}$(>0)$ is not a p-value.} We use it
only to summarize the direction of the bootstrap distribution.
\end{itemize}

% ---------------------------------------------------------------------
% 7. SECTION 9 -- Conclusion (replace whole section)
% ---------------------------------------------------------------------
\section{Conclusion}

We presented an interpretable, leakage-controlled pipeline for early
warning of new negative trust-edge emergence in the Bitcoin-OTC
temporal signed network, together with a three-way feature ablation
under Logistic Regression and Random Forest. We evaluated under a
rolling-origin protocol over 25 disjoint single-month test folds
(425 total positive events), with paired Wilcoxon and sign tests on
per-fold metric deltas and a pre-declared primary family controlled
by Holm-Bonferroni.

\textbf{Temporal features robustly improve early warning.} The
$(\mathrm{S}{+}\mathrm{T}) - \mathrm{S}$ ROC-AUC delta is positive in
22 of 25 RF folds and 21 of 25 LR folds, with Wilcoxon p-values
$\leq 1.4\!\times\!10^{-4}$ under both models after Holm correction
and matched-pairs rank-biserial $r_{\mathrm{rb}} \approx +0.8$. The
Random Forest temporal PR-AUC lift is also significant after Holm
correction ($p = 6.3\!\times\!10^{-4}$,
$r_{\mathrm{rb}} = +0.74$).

\textbf{Community-aware features provide a small, PR-oriented gain
under Random Forest.} The
$(\mathrm{S}{+}\mathrm{T}{+}\mathrm{C}) - (\mathrm{S}{+}\mathrm{T})$
PR-AUC delta is positive in 19 of 25 RF folds (sign-test
$p = 0.015$; Wilcoxon $p = 0.048$, marginal under the pre-declared
Holm threshold; $r_{\mathrm{rb}} = +0.45$). The effect is real
enough to mention but considerably smaller than the temporal effect.

\textbf{Community features do not improve ROC-AUC over the temporal
baseline.} Neither RF nor LR shows a significant community ROC-AUC
gain ($p = 0.096$ and $p = 0.47$ respectively), and the RF
direction is mildly negative on average -- consistent with
saturation.

\textbf{Logistic Regression shows no community gain on any metric},
consistent with the community signal being non-linear.

\textbf{Precision@$k$ is unstable under the rolling protocol} --
most folds return no true positives in the top 50 -- and we
therefore demote it from headline status to a supplementary metric.

At the descriptive level, the data are not consistent with a coarse
``distrust between communities'' reading: same-community pairs in
aggregate training data have a higher empirical positive rate than
cross-community pairs. They are consistent with a finer pair-level
pattern in which lower neighborhood overlap is associated with
higher new-negative probability. We frame this as a refined,
conservative reading rather than as a confirmed weak-tie result.

The largest natural extensions are: (i) stronger baselines,
including recency-only and victim-risk-only heuristics and
gradient-boosted trees; (ii) external replication on Bitcoin-Alpha;
(iii) a richer conditional analysis of community signal (e.g.\
overlap-conditioned positive rates inside and across communities).
